% !PRINTED 321
\weberpage{321}
%<<CONT
gekehrt leicht zeigen, dass die vier in der Form (8) enthaltenen
Grössen $X$, wenn die $v_{1,0}, v_{0,1}, v_{1,1}$ die Wurzeln irgend einer
cubischen Gleichung sind, einer biquadratischen Gleichung mit
rationalen Co\"efficienten genügen\textsuperscript{1)}.

\noindent\rule{0.35\linewidth}{0.4pt}

\footnotesize
\noindent\textsuperscript{1)} Diese Form der Lösung der biquadratischen Gleichung ist das
Analogon zu der \gesperrt{Cayley'schen} Form der \gesperrt{Cardanischen Formel} (Bd. I,
\S.\,36). Ich weiss nicht, ob diese Form der Auflösung biquadratischer
Gleichungen in der elementaren Algebra sonst schon bekannt ist. Ich
habe sie zuerst in einer englischen Aufgaben-Sammlung \glqq Mathematical
Tripos, Part I, June 1894\grqq{} gefunden, wo auch die cubische Resolvente
der $v$ gebildet ist.
\normalsize
